\section{Заключение}

\subsection{Сводные результаты}

В формальной постановке задачи определена целевая функция
$t_{\mathrm{total}}(\mathcal{A}) = t_{\mathrm{plan}}(\mathcal{A}) + t_{\mathrm{exec}}(\mathcal{A}) \to \min_{\mathcal{A}}$
и обозначен конфликт между точностью перебора (малое $t_{\mathrm{exec}}$)
и его стоимостью (большое $t_{\mathrm{plan}}$).
Для решения этого конфликта в работе предложены и реализованы два подхода
к адаптивному планированию запросов в ядре PostgreSQL.
Все подходы протестированы на двух классах запросов $\mathcal{Q}$:
JOB \cite{HowGoodJournal} и JOB-Complex \cite{JOBComplex}.
Результаты нормализованы относительно стандартного DPsize:

\begin{center}
\small
\begin{tabular}{l c c c c | c c c c}
 & \multicolumn{4}{c|}{JOB (113 запросов)} & \multicolumn{4}{c}{JOB-Complex (30 запросов)} \\
 & $t_{\mathrm{plan}}$ & $t_{\mathrm{exec}}$ & $t_{\mathrm{total}}$ & побед
 & $t_{\mathrm{plan}}$ & $t_{\mathrm{exec}}$ & $t_{\mathrm{total}}$ & побед \\
\hline
Связывание якорных участков
    & $\mathbf{0{,}78}$ & $\mathbf{0{,}89}$ & $\mathbf{0{,}85}$ & \textbf{102}
    & $0{,}85$ & $0{,}59$ & $0{,}61$ & \textbf{20} \\
Декомпозиция на топологии
    & $0{,}83$ & $3{,}55$ & $2{,}55$ & 32
    & $0{,}82$ & $3{,}12$ & $2{,}99$ & 12 \\
\end{tabular}
\end{center}

\subsection{Выводы}

В аннотации к работе поставлена цель: получить глубокое понимание того,
как разные алгоритмы планирования возможно комбинировать,
делая переключение между ними в зависимости от топологии и сложности
графа запроса.
Реализация в ядре PostgreSQL и экспериментальная оценка на двух бенчмарках
позволяют сформулировать следующие результаты.
\par

\textbf{1. Замена алгоритма перечисления подпланов при сохранении
пространства поиска не является достаточным средством
минимизации $t_{\mathrm{total}}$.}
DPhyp и его модификация DPhypMod рассматривают то же множество пар, что DPsize,
и на JOB дают нейтральный или отрицательный результат.
Причина: свыше 90\% времени $t_{\mathrm{plan}}$ приходится
на процедуру построения путей, число вызовов которой определяется
топологией графа, а не алгоритмом обнаружения пар.
Для сокращения $t_{\mathrm{plan}}$ необходимо уменьшать
само пространство поиска.
\par

\textbf{2. Результат планирования зависит от порядка перебора,
и эта зависимость усиливается с ростом ошибок оценок.}
Процедура инкрементального отбора путей в PostgreSQL сравнивает
стоимости с допуском $\varepsilon = 1\%$; при совпадении побеждает
путь, поступивший раньше.
DPsize обрабатывает левостороние подпланы первыми, создавая
систематическое смещение.
На JOB это смещение нейтрально.
На JOB-Complex оно приводит к выбору планов с вложенными циклами с миллиардами итераций.
Жадные алгоритмы $\mathrm{GOO_{card}}$ и $\mathrm{GOO_{comb}}$
избегают этих ловушек, поскольку не опираются на стоимостную модель
при сравнении подпланов большого размера, и находят планы
на 1--2 порядка быстрее DPsize.
\par

\textbf{3. Качество оценок кардинальностей определяет $t_{\mathrm{exec}}$,
а не выбор алгоритма перечисления.}
Это основной вывод работы.
На JOB, где ошибки оценок умеренны, полный перебор DPsize
даёт хорошие планы, и якорное связывание улучшает его.
На JOB-Complex, где ошибки достигают $10^5$--$10^8$,
DPsize выбирает планы в $11\times$ медленнее оптимальных \cite{JOBComplex},
и $\mathrm{GOO_{card}}$~--- жадный алгоритм без стоимостной модели~---
оказывается в $9\times$ лучше DPsize.
Стоимостная модель PostgreSQL, построенная на предположении
независимости атрибутов, при нарушении этого предположения
не просто теряет точность, а \emph{активно вредит}, направляя
перебор в область катастрофических планов.
\par

\subsection{Направления дальнейшей работы}

\textbf{1. Динамический подбор параметров для якорного подхода.}
Подбор параметров во время планирования позволил бы сделать данный 
подход более гибким к разным классам запросов, схемам баз данных.
\par

\textbf{2. Обнаружение ненадёжных оценок.}
PostgreSQL не предоставляет индикатора уверенности
в оценке кардинальности.
Разработка лёгковесной эвристики, маркирующей соединения
с потенциально ненадёжными оценками
(строковые столбцы, отсутствие первичного или внешнего ключа в соединении, 
высокий $n_{\mathrm{distinct}}$),
позволила бы использовать кардинальностную метрику, а для 
надёжных~--- стоимостную,
приближая реализацию к алгоритму $\mathcal{A}^*$,
адаптивному не только по сложности графа,
но и по качеству доступных статистик.
\par

3. Использование стоимостных критериев и селективности для определения якорей и их участков 
на единичных отношениях.
\par

4. Выработка других критериев для применения якорного подхода.